\subsection{Automorphisms of Semisimple Algebras}

THEOREM 4. --- Let A be a semisimple ring, Z its center, and u an automorphism of A. Suppose that A is a finitely generated Z-module and that we have \(u(z) = z\) for every z in Z. Then u is an inner automorphism.

This follows from Theorem 2 of VIII, p.~256 applied with \(f = \mathrm{Id}_{\mathrm{A}}\) and \(g = u\) .

Example. --- Theorem 4 applies in the following two specific cases:

\pandocbounded{\includegraphics[keepaspectratio]{images/part002_1d1a6ca2.jpg}}

\begin{enumerate}
\def\labelenumi{\alph{enumi})}
\item
  Let D be a field and Z its center. If D has finite degree over Z, then every automorphism of D that fixes the elements of Z is an inner automorphism. The assumption that D has finite degree over Z is essential (VIII, p.~269, Exercise 4).
\item
  Let V be a finite-dimensional vector space over the field K. Every automorphism of the K-algebra \(\operatorname{End}_{\mathrm{K}}(\mathrm{V})\) is an inner automorphism; this result extends to the case when the space V is not finite-dimensional over K (VIII, p.~272, Exercise 13).
\end{enumerate}

In particular, every automorphism of a matrix algebra \(\mathbf{M}_{n}(\mathrm{K})\) (with \(n \geqslant 1\) ) is an inner automorphism. This result admits the following generalization.

PROPOSITION 2. --- Let L be a commutative ring and V a free L-module of finite dimension m. Suppose that every L-module M such that \(M^{m}\) is isomorphic to \(L^{m}\) is isomorphic to L. Then every automorphism of the L-algebra \(\operatorname{End}_{\mathrm{L}}(\mathrm{V})\) is an inner automorphism.

Set \(B = \text{End}_{\text{L}}(\text{V})\) . Let u be an automorphism of the L-algebra B. We view V as a left B-module; let \(u_{*}(\text{V})\) be the left B-module associated with u, with external law \((b, v) \mapsto u(b)(v)\) (II, §1, No.~13, p.~221). Let \((e_1, \ldots, e_m)\) be a basis of the L-module V; given elements \(v_1, \ldots, v_m\) of V, there exists a unique element b of B such that we have \(b(e_i) = v_i\) for \(1 \leqslant i \leqslant m\) . In other words, the element \(e = (e_1, \ldots, e_m)\) of \(V^m\) provides a basis of the B-module \(V^m\) . Since u is an automorphism, e also gives a basis of the B-module \(u_{*}(\text{V}^m) = u_{*}(\text{V})^m\) , which is therefore isomorphic to \(V^m\) . The (B,L)-bimodule V is invertible (VIII, p.~102, Example 1). By Theorem 2, b) of VIII, p.~103, there consequently exists an L-module M such that the B-module \(u_{*}(\text{V})\) is isomorphic to \(V \otimes_{L} M\) . The B-modules \(V \otimes_{L} L^m\) and \(V \otimes_{L} M^m\) , respectively isomorphic to \(V^m\) and \(u_{*}(\text{V})^m\) , are therefore isomorphic. By loc. cit., the L-modules \(L^m\) and \(M^m\) are isomorphic. Given the assumption, M is isomorphic to L; consequently, the B-module \(u_{*}(\text{V})\) , which is isomorphic to \(V \otimes_{L} M\) , is isomorphic to V. Let h be a B-module isomorphism from V to \(u_{*}(\text{V})\) ; it is, in particular, an automorphism of the L-module V, that is, an invertible element of B. For b in B and v in V, we have \(h(b(v)) = u(b)(h(v))\) and therefore \(u(b) = hbh^{-1}\) .

The conditions of Proposition 2 are, in particular, satisfied when the commutative ring L is a principal ideal domain (VII, §3, p.~15, Corollary 3) or is Artinian (VIII, p.~37, Theorem 2, d)) or local (VIII, p.~36, Corollary 6).

